\section{Literal contours and finite sets of states}
\label{geo:section}

This section constructs many families of cyclic bases that remain disjoint
after certain local changes. The local changes will later prevent
intersections between families built on different sets of control symbols.
The construction keeps track of the actual order of the symbols. This is
needed because the word assembly uses literal overlaps, not just cyclic
equivalence classes.

Fix an integer $r\geq 6$, and write $k=r-1$ and $d=r-3$. There are $M$
distinct \emph{background symbols}, in a fixed cyclic order, and $r$
distinct \emph{control symbols}. These two alphabets are disjoint. Number
the background symbols by the elements of $\mathbb Z/M\mathbb Z$, in their
cyclic order. A \emph{cut} is the position immediately before a background
symbol; its coordinate is the number of that symbol. There are $k$
stationary positions for control symbols and one mobile position. The
stationary positions have fixed identities $0,\ldots,k-1$. Their identities
do not change when their cuts pass the distinguished background symbol.

A \emph{literal base} is a list containing all $M+r$ symbols exactly once.
Two bases have the same \emph{cyclic class} if one is a cyclic rotation of
the other. We always write the mobile symbol last. After deleting it, the
remaining list is a representative of the stationary cyclic order. Thus
the stationary cuts and their resident control symbols, together with the
gap occupied by the mobile symbol, determine the literal base uniquely.
The gap is necessary: the stationary cuts and the assignment of control
symbols alone do not determine a base during a sequence of full moves.

Put $h=M+k$, so a base has length $h+1$. For a base
$b=b_0\cdots b_h$, define two literal operations by
\[
 F(b)=b_1\cdots b_{h-1}b_0b_h,
 \qquad
 S(b)=b_hb_0\cdots b_{h-3}b_{h-1}b_{h-2}.
\]
A \emph{full row} consists of a base with visible length $h+1$, and a
\emph{short row} consists of a base with visible length $h-1$. Their
\emph{charges} are respectively $0$ and $2$. A fixed extra symbol, called
the satellite, is outside the base alphabet and is attached to every row
when the rows are used in the word construction. It does not change in
this section.

The head of a row is the first $h-1$ symbols of its base. Its tail is the
first $h-1$ symbols after rotating the base left by one more than its
visible length. Two consecutive rows are compatible when the first tail
equals the second head. The displayed formulas show directly that the
successor of a full row with base $b$ may have base $F(b)$, and the
successor of a short row may have base $S(b)$. A \emph{closed contour} is a
nonempty cyclic list of compatible rows. Its source classes are the cyclic
classes of its bases, with one occurrence for each row.

\subsection{Positive updates and local simplicity}

Suppose a stationary position contains $A$, the background symbol at its
cut is $x$, and the mobile symbol is $T$. Apply $F$ until the base ends in
$A,x,T$, and then apply $S$. We call this a \emph{positive update} of that
stationary position. The cyclic order changes locally from
\[
 \cdots A,x,T\cdots
 \quad\hbox{to}\quad
 \cdots x,A,T\cdots.
\]
Consequently the old mobile symbol $T$ becomes the stationary resident at
the next background cut, and $A$ becomes mobile. The new mobile is
immediately before the new resident. In terms of the assignment of control
symbols to positions, the update exchanges the mobile position with the
chosen stationary position. No other stationary cut changes.

A \emph{batch} updates every stationary position once, in some prescribed
order. If the initial cuts are $q_0,\ldots,q_{k-1}$, then during batch $t$
the current cuts are
\begin{equation}
 q_i+t+\epsilon_i\pmod M,
 \qquad \epsilon_i\in\{0,1\}.
 \label{geo:cuts}
\end{equation}
Here $\epsilon_i$ is one exactly when position $i$ has already been
updated in the current batch. After the last update, increase $t$ by one
and reset all the $\epsilon_i$ to zero. The short row used for that last
update still has the old value of $t$ and has phase zero at its last
target. Only its successor uses the reset convention. Full moves change
the mobile gap and leave the data in \eqref{geo:cuts} unchanged.

The circular distance between two residues modulo $M$ is the smaller of
the two distances between them around the $M$ cuts. Call the initial cuts
\emph{robust} if the indexed family
\[
 0,\qquad q_i-q_j\quad(0\leq i,j<k,\ i\ne j)
\]
has circular distance at least five between any two different members.
The word ``indexed'' matters: repeated differences are forbidden. In
particular, the cuts are distinct. Perturbing each cut by zero or one
leaves distinct stationary cuts at circular distance at least four.

For example, with $r=6$ and $M=360$, the cuts
\[
 (q_0,q_1,q_2,q_3,q_4)=(0,5,25,125,155)
\]
are robust. Their positive integer differences are
$5,20,25,30,100,120,125,130,150,155$. These differences, their negatives
modulo $360$, and zero are separated by at least five.

\begin{lemma}
\label{geo:proper-run}
Assume the cuts are robust. After a positive update, the next positive
update reaches its required suffix in fewer than $h-1$ full moves.
\end{lemma}
\begin{proof}
The stationary list has $h$ distinct symbols. Each full move rotates it
one position, so the required target suffix is reached in a unique number
$u$ of moves with $0\leq u<h$. It remains to exclude $u=h-1$.

Let $x$ be the background symbol crossed by the preceding short move.
Immediately after that move, the stationary list ends in $x$, and the
position just updated has cut $x+1$. If $u=h-1$, its final symbol $x$
becomes its first symbol. On the other hand, let $z$ be the cut of the
next target. At the required suffix, the stationary list ends in that
target's resident followed by $z$. The successor cut $z+1$ is unoccupied,
because the current stationary cuts are separated by at least four.
Therefore the first stationary symbol at this suffix is the background
symbol $z+1$. We obtain $x=z+1$ modulo $M$.

The previous target now has cut $x+1=z+2$, whereas the next target has
cut $z$. If these are different stationary positions, their distance two
contradicts the separation just stated. If they are the same position,
the equality would give $2=0$ modulo $M$, which is impossible: robust
cuts include at least two different cuts at circular distance at least
five, so $M\geq10$. Thus $u\ne h-1$, proving the stronger bound.
\end{proof}

\begin{lemma}
\label{geo:simplicity}
Assume the cuts are robust. Consider a schedule of exactly $M$ batches
that returns the complete literal base to its initial value. Its source
classes are pairwise distinct.
\end{lemma}
\begin{proof}
Suppose two emitted rows have the same cyclic base. If their mobile
symbols are equal, they have $k$ common stationary control symbols;
otherwise they have $k-1$. In either case there are at least two. A common
stationary symbol occupies some position $i$ in the first row and some
position $j$ in the second. Equality of its background cut gives
\[
 t'-t=q_i-q_j+\epsilon_i-\epsilon'_j\pmod M.
\]
The ordered pairs $(i,j)$ for different common symbols are distinct. Two
off-diagonal pairs would place distinct nonzero ordered differences within
distance two. One diagonal and one off-diagonal pair would place a
nonzero difference within distance two of zero. Robustness excludes both
possibilities. Therefore every common symbol occupies the same stationary
position in both rows, and $t'-t$ is $0$, $1$, or $-1$ modulo $M$.

Choose the direction of comparison so that the rows lie in the same batch
or in consecutive batches of the cyclic schedule. Literal closure permits
this comparison across the last and first batches. In the same batch,
equality of the phases of a common stationary symbol implies that its
position is not updated between the rows. In consecutive batches, its
phase must be one in the first row and zero in the second. It has already
been updated before the first row, and is not updated again before the
second. Thus no common stationary position is updated in the interval
between the rows.

If the mobile symbols are equal, every stationary position is common.
The two rows lie in a single run of full moves, possibly including its
final short source row. The $h$ possible insertion gaps for the mobile
in the fixed stationary circle give $h$ different cyclic bases. By
Lemma~\ref{geo:proper-run}, that run has fewer than $h$ full moves, so no
two of its source rows have the same cyclic base.

If the mobile symbols differ, write them as $a$ and $b$. The only
stationary position not common to the two rows is one position $j$,
occupied by $b$ first and by $a$ second. All intervening updates must be
at $j$. There is at most one such update in one batch, and at most two
across the boundary between consecutive batches. The mobile changes, so
the number of these exchanges is odd. It is therefore exactly one.

Let $x$ be the cut of $j$ before that update. Before the short move, $b$
is stationary at $x$ and $a$ is mobile. Afterwards, $a$ is stationary at
$x+1$ and $b$ is mobile. Equality of the two cyclic bases requires $a$
to occupy cut $x+1$ before the update. There is only one such mobile gap,
because no stationary resident is at that cut before the update. This is
the gap of the short source row. The same equality requires the new
mobile $b$ to occupy cut $x$ afterwards. Immediately after the short move,
the stationary list begins with $a$ and ends with the background symbol
$x$. Inserting $b$ before $x$ therefore requires the final full gap,
reached after exactly $h-1$ full moves. Lemma~\ref{geo:proper-run} says
that the next target is reached strictly before this gap. The required
row therefore does not occur. This excludes the final possibility.
\end{proof}

\subsection{The normal order and exact replacement intervals}

Choose representatives $0=q_0<q_1<\cdots<q_{k-1}<M$. The \emph{normal
order} is
\[
 O=(k-1,k-2,\ldots,0).
\]
Begin with the mobile immediately before the resident at position zero.
As before, the mobile is the last symbol of the literal base. A complete
normal batch acts on the $r$ control positions as an $r$-cycle, denoted
$C$. Indeed, the original mobile goes to the first target; each subsequent
target receives the resident displaced from the preceding target; and the
resident of the last target becomes mobile.

\begin{lemma}
\label{geo:normal-counts}
If $r$ divides $M$, then $M$ normal batches form a literally closed
contour. Its number of rows $N_0$ and its charge $Q_0$ are
\[
 N_0=M\bigl((r-2)M+r(r-1)+1\bigr),
 \qquad Q_0=2(r-1)M.
\]
\end{lemma}
\begin{proof}
The role assignment returns because $C^M$ is the identity. Every
stationary cut advances $M$ times and therefore returns. The final target
is again zero, so the mobile is again immediately before its resident.
These data determine the same literal base, with the mobile last.

In one normal batch the first full run crosses $q_{k-1}+1$ background
symbols and all $k$ resident controls. After updating target $i$, the
full run to target $i-1$ crosses $M+q_{i-1}-q_i$ background symbols and
all $k$ resident controls. Consequently the total number of full moves is
\[
 q_{k-1}+1+k+
 \sum_{i=1}^{k-1}(M+q_{i-1}-q_i+k)
 =(k-1)M+1+k^2.
\]
There are also $k$ short rows, each of charge two. Multiplying the two
batch totals by $M$ gives the formulas.
\end{proof}

We next describe a replacement that keeps a specified control symbol out
of a specified stationary position for a chosen interval. Let $P$ be an
order of the $k$ stationary positions, with first target $a$. Performing
$P$ followed by its reverse is a \emph{parking pair}. Each position is
updated twice. The exchanges of roles cancel in reverse order, so every
control returns to its original role. The control that was initially
mobile is placed at $a$ by the first exchange and recovered by the last.
No intermediate exchange uses $a$. Throughout the pair, including all
full rows, this control is either mobile or resident at $a$.

For example, if $k=5$ and the forbidden position is $4$, use
$P=(3,4,2,1,0)$. A control initially mobile stays out of position $4$
throughout this order and its reverse. At the midpoint it is resident at
position $3$, although all stationary cuts have already advanced once.

\begin{lemma}
\label{geo:parking-endpoint}
Let $T$ be a positive integer. Replace $2rT+1$ consecutive normal batches
by $rT$ parking pairs, all with the same order $P$, followed by one normal
batch. The replacement has exactly the same final literal base as the
normal interval and the same total number of short rows. A control
initially mobile never occupies any stationary position different from
the first target of $P$ during the parking portion.
\end{lemma}
\begin{proof}
The parking pairs restore the role assignment, and the final normal
batch applies $C$. The normal interval also applies $C$, because
$C^{2rT+1}=C$. Both intervals advance every stationary cut $2rT+1$ times.
Both end with an update of position zero. The post-update mobile gap is
therefore the same, and the literal bases agree. Each interval has
$k(2rT+1)$ short updates. The final assertion follows from the preceding
description of each pair. Protection is asserted during the parking
portion; the final normal batch must lie outside the interval to be
protected.
\end{proof}

We call the final normal batch the \emph{cleanup batch}. It is needed
because a parking pair can end with its mobile immediately before a
resident other than that of position zero. Equality of role assignments
alone would not give equality of literal endpoints.

\subsection{A catalogue of joint role and phase states}

Only two parking orders are needed. Put $P_0=O$, and obtain $P_1$ by
exchanging the first two entries of $O$. For a forbidden position $j$,
use $P_1$ when $j$ is the first target of $P_0$, and use $P_0$ otherwise.
The chosen first target always differs from $j$. Fix this choice for all
parking pairs in a replacement interval.

An \emph{admissible schedule} is obtained from the normal $M$-batch
contour by replacing mutually disjoint intervals as in
Lemma~\ref{geo:parking-endpoint}, using only $P_0$ and $P_1$. An interval
may be placed on the cyclic schedule; choose the initial cut outside the
replacement intervals. Since each replacement has its normal literal
endpoint, all replacements can be performed without changing the states
at the other normal boundaries.

Let a role assignment be a bijection from the $r$ positions, including
the mobile position, to the $r$ named controls. Write $U_\ell(P)$ for the
role action of the first $\ell$ updates of $P$ followed by its reverse,
where $0\leq\ell\leq 2k$. For such a prefix define a phase vector by
\[
 \epsilon_\ell(P)=
 \begin{cases}
  \text{the indicator of the first $\ell$ targets of $P$},&0\leq\ell<k,\\
  \text{the indicator of the first $\ell-k$ targets of $P^{\rm rev}$},
       &k\leq\ell<2k,\\
  0,&\ell=2k.
 \end{cases}
\]
At the midpoint $\ell=k$, the translation has increased by one and the
phase vector is zero. The role action remains $U_k(P)$; it has not reset.
At $\ell=2k$, the translation has increased by two and both the phase
vector and the role action have reset.

Define the finite set of \emph{joint states}
\begin{equation}
 \mathcal J=
 \left\{\bigl(C^sU_\ell(P),\epsilon_\ell(P)\bigr):
      0\leq s<r,\ P\in\{P_0,P_1\},\ 0\leq\ell\leq2k\right\}.
 \label{geo:joint-set}
\end{equation}
Composition here means chronological action on a role assignment: the
normal batches act first, followed by the indicated prefix. We count
joint states, not independent choices of a role permutation and a phase
vector.

\begin{lemma}
\label{geo:joint-containment}
Fix an initial role assignment $\sigma$. Every source row of an admissible
schedule has stationary cuts $q_i+t+\epsilon_i$ and role assignment
$\sigma\tau$ for some translation $t$ and some
$(\tau,\epsilon)\in\mathcal J$. Moreover,
\[
 |\mathcal J|\leq2r(2k+1)=4r^2-2r<4r^2.
\]
\end{lemma}
\begin{proof}
A normal boundary has assignment $\sigma C^s$ for some $s$ modulo $r$.
The proper prefixes of its batch are first-half prefixes of the $P_0$
pair. Its final state is the empty prefix at normal phase $s+1$.

A replacement starts at a normal boundary. Every complete parking pair
restores its starting role assignment. Thus every pair in the replacement
has the same list of relative role actions $C^sU_\ell(P)$. The reset
convention in \eqref{geo:cuts} gives exactly the displayed phase vectors,
including the midpoint. The cleanup has the proper prefixes of a $P_0$
batch, and its final state is the next normal role phase. By
Lemma~\ref{geo:parking-endpoint}, the subsequent normal boundary agrees
with the original schedule, so the argument applies successively to all
replacements.

Full moves change neither the role assignment nor the stationary cuts.
The short row preceding the last update of a batch is represented by the
proper prefix before that update; its successor is represented at the
reset boundary. Thus every emitted row has been included. Finally, $r$
divides $M$, the translation increases by $M$, and the schedule closes
literally. The last and first states obey the same convention. The
cardinality bound counts the indices in \eqref{geo:joint-set}; repeated
representations only decrease the number of states.
\end{proof}

\subsection{Many robust modular frames}

The stationary cut set will have one additional algebraic property that
allows us to reconstruct a missing cut. Choose integers
\[
 0<z_1<\cdots<z_d<M,
 \qquad s=z_1+\cdots+z_d\pmod M,
\]
where $s$ is represented in $\{0,\ldots,M-1\}$. Sort the $k=d+2$ points
$0,z_1,\ldots,z_d,s$ into increasing order, and retain the result only
when it is robust. This ordered cut tuple is a \emph{modular frame}.
The point $s$ is its sum point. Assign coefficient $d-1$ to zero,
coefficient $1$ to the sum point, and coefficient $-1$ to every $z_i$.
These coefficients are all nonzero, their sum is zero, and their weighted
sum of cuts is zero modulo $M$. Thus the relation is unchanged by a
common translation. In the example above, $z=(5,25,125)$ and $s=155$.

We use the following elementary fact for arbitrary moduli. For a nonzero
integer $a$, the congruence $ax=b\pmod M$ has either no solutions or
$\gcd(|a|,M)$ solutions in $\mathbb Z/M\mathbb Z$. Indeed, if one solution
exists then all solutions form a translate of the kernel of multiplication
by $a$. On writing $g=\gcd(|a|,M)$ and dividing by $g$, coprimality shows
that the kernel consists of the $g$ multiples of $M/g$. In particular,
there are at most $|a|$ solutions. No primality assumption is needed.

\begin{lemma}
\label{geo:frame-supply}
Let $F_M$ be the number of robust modular frames. Then
\[
 F_M\geq
 \frac{(M-1)^d-18k^4(M-1)^{d-1}}{d!}
 \geq\frac{M^d}{d!}\left(1-\frac{19r^4}{M}\right).
\]
Each retained anchored cut set has a unique sum point.
\end{lemma}
\begin{proof}
First choose an ordered vector from $\{1,\ldots,M-1\}^d$, without sorting
its coordinates. Regard the cuts as evaluations of the formal points
$0,e_1,\ldots,e_d,\boldsymbol 1$ in $\mathbb Z^d$. Their affine relations
form a one-dimensional space, generated by
\[
 (d-1,-1,\ldots,-1,1).
\]
To verify this, a relation with coefficient $a$ at $\boldsymbol 1$ has
coefficient $-a$ at each $e_i$, by comparing coordinates, and coefficient
$(d-1)a$ at zero, by summing coefficients. Every nonzero relation
therefore involves all $d+2\geq5$ points.

A comparison of two distinct ordered differences, or of a nonzero
difference with zero, involves at most four formal points. Its linear
form is therefore nonzero. Every coefficient in that form has absolute
value at most two. There are at most $k^4$ comparisons. For robustness,
exclude the nine residues $-4,-3,\ldots,4$ for each comparison. Fix all
but one coordinate with nonzero coefficient. The modular root bound gives
at most two choices for the remaining coordinate for each excluded
residue. The number of bad ordered vectors is at most
$18k^4(M-1)^{d-1}$.

The good vectors are invariant under coordinate permutation, and a good
vector has distinct coordinates. Its orbit under coordinate permutation
therefore has exactly $d!$ elements. Sorting gives the first inequality.
For the second, use
$(M-1)^d\geq M^d-dM^{d-1}$ and
$d+18k^4\leq19r^4$.

It remains to check that sorting has not identified different sum
decorations. If an anchored cut set had two sum points $s$ and $s'$, the
sum of all its cuts would be both $2s$ and $2s'$ modulo $M$. Distinct
solutions require $M$ even and $s-s'=M/2$. The two ordered differences
$s-s'$ and $s'-s$ would then coincide. This contradicts robustness, so
the sum point is unique.
\end{proof}

\subsection{Reconstruction and disjoint labelled families}

A \emph{labelled frame} is a pair $(q,\sigma)$, where $q$ is a robust
modular frame and $\sigma$ is an initial bijective assignment of the $r$
named controls to the $r$ positions. Its \emph{envelope} is the set of
cyclic bases obtained by choosing an arbitrary translation, a joint state
$(\tau,\epsilon)\in\mathcal J$, and an arbitrary mobile insertion gap,
with stationary cuts $q_i+t+\epsilon_i$ and assignment $\sigma\tau$.
Lemma~\ref{geo:joint-containment} shows that every source row of every
admissible repair of this labelled frame lies in its envelope.

\begin{lemma}
\label{geo:reconstruction}
Fix a labelled frame and fix a joint state on each side of a possible
envelope intersection. The number of positional reconstruction witnesses
for a target frame is at most
\[
 D_0=k(k-1)^2\bigl(d+1+\gcd(d-1,M)\bigr)
       \leq2d\,k(k-1)^2\leq2r^4.
\]
A witness includes the matching of the retained stationary positions.
\end{lemma}
\begin{proof}
A common cyclic base has at least $k-1$ controls stationary on both sides.
If the mobile symbols agree, discard one of the $k$ common stationary
controls. Choose the omitted stationary position on the source side and
on the target side. The remaining $k-1$ stationary controls occur in the
same cyclic order. The binary phase perturbations do not change their
order, because the cuts remain separated by at least four. There are
therefore $k-1$ possible cyclic matchings.

For a matched source position $i$ and target position $j$, equality of
the cuts gives
\[
 q'_j+\delta=q_i+\epsilon_i-\epsilon'_j\pmod M,
\]
where $\delta$ is the difference of the two common translations. The
phase vectors are fixed, so the right-hand side is known. Choose also
the target sum position $a\in\{1,\ldots,k-1\}$. This fixes its coefficient
vector $c$. Set $w_j=q'_j+\delta$. Since $\sum c_j=0$, the target relation
is $\sum c_jw_j=0$. It determines the omitted coordinate $w_j$ in at most
$\gcd(|c_j|,M)$ ways. If position zero was retained, its known coordinate
$w_0$ determines $\delta$; if it was omitted, its reconstructed coordinate
does so. Subtracting $\delta$ then determines every target cut. Sorting
and robustness can reject a reconstruction but cannot add one.

For each fixed sum position, summing the root multiplicities over the
omitted target position gives
$d+1+\gcd(d-1,M)\leq2d$. There are $k$ source omissions, $k-1$ cyclic
matchings, and $k-1$ choices of sum position. Multiplying gives the stated
bound, which counts witnesses even if several reconstruct the same frame.
\end{proof}

\begin{lemma}
\label{geo:labelled-packing}
Put
\[
 \Delta=64r^8,
 \qquad \eta=\frac{r!}{2d!(\Delta+1)}.
\]
For every $M\geq38r^4$, there are
$\lfloor\eta M^d\rfloor$ labelled frames whose envelopes are pairwise
disjoint.
\end{lemma}
\begin{proof}
Form a finite graph whose vertices are all labelled frames, with an edge
when two distinct vertices have intersecting envelopes. Distinct vertices
with coincident envelopes are also adjacent. By
Lemma~\ref{geo:frame-supply}, the graph has at least
$r!M^d/(2d!)$ vertices.

Fix a source vertex. For a chosen pair of joint states, each witness of
Lemma~\ref{geo:reconstruction} determines the target frame and matches
$r-2$ named stationary control symbols. Since the target role action is
a permutation, those equalities determine the target initial assignment
on $r-2$ positions. The remaining two labels have at most two assignments
to the remaining positions. Thus the number of possible neighbouring
vertices, including the source vertex itself, is at most
\[
 2|\mathcal J|^2D_0\leq2(4r^2)^2(2r^4)=64r^8.
\]
This bound includes all assignments at the time the graph is formed. It
would not be valid to select unlabelled frames first and then attach all
$r!$ assignments to each selected frame.

Greedily select a vertex and remove it and all its neighbours. Each
selection removes at most $\Delta+1$ vertices. Continuing until no
vertices remain produces an independent set with at least
$r!M^d/(2d!(\Delta+1))=\eta M^d$ vertices. Select the required integer
number from this independent set. Their envelopes are disjoint by the
definition of the graph.
\end{proof}

For fixed $M$, the selected family can be used for every background circle
and every set of $r$ active letters. To see this, first make the selection
using $M$ formal background symbols and $r$ different formal control
symbols. Replace the background symbols, in order, by the letters of any
chosen background circle, and replace the formal controls by its active
letters. These substitutions give an injective map on the entire formal
alphabet. They commute with the construction of each row and with cyclic
rotation. If two substituted rows were equal up to rotation, injectivity
would give the same equality before substitution. Thus the selected
envelopes remain disjoint. In particular, the choice of frames is
independent of the guard positions introduced in the next section.

\subsection{What the local construction supplies to assembly}

For later use, the literal displacement also has an exact description.
Give each background symbol its index modulo $h$ in a literal base. A
full move changes that index by $-1$. A short move changes it by $+1$,
except for the crossed background symbol, which stays at index $h-1$.
In $M$ complete batches each stationary position crosses every background
symbol exactly once. Every background symbol therefore has exactly $k$
exceptional short moves. If the contour is literally closed, with $N$
rows and total charge $Q$, this proves
\begin{equation}
 Q-N\equiv k\pmod h,
 \qquad \gcd(|Q-N|,h)=\gcd(k,h)\leq k.
 \label{geo:displacement}
\end{equation}

An untouched normal batch contains an actual consecutive run of literal
$F$ successors of length at least
\[
 \frac{(k-1)M+k^2+1}{k}.
\]
It is an untouched batch that gives this estimate. The first full run of
a cleanup batch can start at a different mobile gap and is not covered
by the normal count. The global repair argument must therefore leave an
untouched batch before applying a word assembly lemma that requires such
a run.

The results in this section construct disjoint envelopes for one active
control alphabet on one fixed background. They also prove that every
admissible repair stays inside its envelope. To use several overlapping
active alphabets, one must additionally construct the forbidden-event
catalogue, prove the deletion estimates, and place disjoint replacement
intervals. An actual word bound further requires the lists of unused
source classes, an iteration that inserts a variable number of controls,
and additional cyclic words that meet every required literal overlap.
Those are separate mathematical statements; the numerical counts in this
section do not assume their existence.
